% Fragmento conservado; véase MANIFIESTO_ANTECEDENTES_LECTURA.json.
% Adición propuesta al inicio de registro_k.tex, después de su introducción.
% Conserva todas las definiciones, fórmulas y demostraciones actuales.
% Procedencia: registro_k.tex, §§k-hadamard/k-transversal/k-kappas;
% k_reversibilidad.tex; incidencia_articulacion.tex, inc-ampl-cuaterna.
\paragraph{Coordonnées réversibles et détermination incidencielle.}
\label{ampl-alfa-k-intro-reversible}
Le résultat qui organise ce chapitre est l'équivalence entre trois
systèmes de coordonnées d'un registre préalablement produit par
APP--TRIT--TPK. La transformation de Hadamard relie le registre
signé au registre entier ; l'extracteur de différences relie
ce dernier à ses variations transversales et à sa composante uniforme :
\[
 U\underset{\mathcal H_{12}}{\overset{\mathcal H_{12}/4}{\longleftrightarrow}}K
 \overset{(D_3,D_4,Q)}{\longleftrightarrow}
 (D_3K,D_4K,Q(K)).
\]
La première équivalence s'établit sur le sous-réseau entier
$\mathcal L_H$ ; la seconde, sur l'image de l'extracteur. Une fois fixés la
base $1000$, les $12$ positions et l'origine de phase, l'évaluation
$K\mapsto\kappa_{\mathrm{per}}$ ajoute un quatrième système de coordonnées
réversible : la multiplication par $1000^{12}-1$ suivie des divisions
euclidiennes restitue le registre complet. Les démonstrations de
\S\ref{subsec:k-hadamard}, \S\ref{subsec:k-transversal} et
\S\ref{subsec:k-kappas} établissent ces équivalences dans leurs domaines
respectifs.

Deux compositions agissent ensuite sur ce registre. La composition
positionnelle combine ses composantes avec les coordonnées de clôture,
de propagation et d'autoéchelle, et sa normalisation détermine $\alpha_A$. La
composition incidencielle applique le seuil cubique et obtient
\[
 K\longmapsto B_K=\{j:K_j\geq729\}=\{4,6,9,10\}.
\]
Le registre intégral permet de reconstruire les deux compositions. Le support
$B_K$ conserve spécifiquement l'appartenance à la classe sélectionnée
par le seuil. L'identité $B_K=H_e\cap P_\pi$ de
(8.12) de l'Article I (provenance du prolongement exceptionnel, conservée dans le registre technique) relie cette extraction entière à l'incidence
des codes régionaux. Le support négatif de son bilan avec $K$
complète le drapeau $B_K\subset H^-_\alpha$ développé dans le
prolongement exceptionnel. Ainsi, la fonction de $K$ est définie par
une transformation inversible et par des applications ultérieures
précises ; l'information supplémentaire de profondeur et de transport demeure
dans l'état dont procède le registre.

